\paragraph{Carry the legal cash explanations through the same internal pairs.}
The period summary can also retain the event-level explanations of each legal
book. \texttt{GroupIndirectCash/1} recomputes the group report and each supplied
\texttt{DeclaredIndirectCash/2} result. Each legal bridge must select exactly the
opening, closing, period and cutoff used by its group member. Its complete
financial-report request and result must equal that member's retained report.
A valid explanation for another company, chart or reporting date therefore
cannot substitute for the required history. Both legal explanations remain in
the result alongside the unchanged group report.

Simply adding the two operating-cash explanations would still retain any
internal payment or receipt classified as operating. For the same period and
currency, write $P_j,W_j,N_j,R_j$ for the profit and adjustments already defined
for legal book $j$. Let $T$ be minus the signed sum of the operating cash flows
in the group's explicitly removed pairs. Each receipt is positive and each
payment negative. Summing the checked legal identities and applying that same
removal gives group operating cash
\[
 C_o^{\mathrm{group}}
 =\sum_j\bigl(P_j+W_j+N_j+R_j\bigr)+T.
\]
Here the supported group method leaves profit unchanged: summed legal profit
equals the group's reported profit. A method that eliminates intragroup profit
would require a further derivation and contract. The adjustment $T$ is calculated
from the named gross cash records; it is never a balancing remainder.

The earlier two classifications both give $T=0$. When both loan flows are
operating, the payment of 30 and receipt of 30 cancel in this sum, although both
still need removal from gross receipts and payments. When the payer records
investing and the payee financing, neither flow enters operating cash. To test
the correction itself, consider two additional supplied classifications. If P
records the payment as operating while S records the receipt as financing,
summed legal operating cash is $-30+20=-10$. Removing the internal payment
gives $T=30$, so group operating cash is 20. If P instead records investing
and S operating, summed operating cash is $30+20=50$ and $T=-30$, again leaving
20. These examples test category composition; they do not determine which
classification is applicable to an actual loan.

The constructed event programs reproduce the previously supplied opening and
closing books, including their memory and all balances. The adapter makes two
explicit representational changes: it translates each old literal cash-flow
identity into the execution engine's hash of event and local-flow identity,
and adds the empty input schema required by a method with no memory writes.
The corresponding pair references and report-request hashes change with those
inputs. After these declared changes, the complete legal requests and results
match, rather than only their cash totals. Both the original snapshot records
and the translation are retained, so the correspondence does not rewrite the
old provenance.

Eight cases cover two charts and the four category combinations. Changed
dependencies, unrelated or forged legal traces, missing owners, stale pair
records and repeated removal are refused. A zero correction still retains both
gross sides of its pair. If the caller omits a pair, the report keeps those
flows and explains the resulting conditional amount; it does not infer a match
from equal values. Accounting, source, control, classification and completeness
admission remain false. Wider consolidation methods and complete notes remain
separate requirements.
